% ==============================================================================
% Plantilla Institucional de Trabajo Terminal — UPIIZ IPN
% Archivo: config/datos.tex
% ==============================================================================
% ESTE ES EL ÚNICO ARCHIVO QUE DEBES EDITAR PARA CONFIGURAR TUS DATOS PERSONALES
% Y LOS DEL PROYECTO. NO ES NECESARIO MODIFICAR main.tex.
% ==============================================================================

% ------------------------------------------------------------------------------
% 1. MODALIDAD DEL TRABAJO TERMINAL
% ------------------------------------------------------------------------------
% Descomenta EXACTAMENTE UNA de las siguientes dos líneas:
\TTItrue   % Para Trabajo Terminal I (Diseño detallado)
%\TTIfalse  % Para Trabajo Terminal II / Reporte Final (Implementación y Validación)

% ------------------------------------------------------------------------------
% 2. INFORMACIÓN DEL PROYECTO
% ------------------------------------------------------------------------------
\newcommand{\tituloProyecto}{Título Completo del Trabajo Terminal en Español}
\newcommand{\tituloIngles}{Full Title of the Final Degree Project in English}

% Línea de trabajo institucional (Elige la que corresponda a tu proyecto):
% I.   Diseño e implementación de un sistema robótico, dispositivos o sistemas mecatrónicos.
% II.  Diseño e implementación de una máquina o mecanismo.
% III. Diseño e implementación de componentes o sistemas electrónicos.
% IV.  Diseño e implementación de sistemas o técnicas de control.
% V.   Diseño y desarrollo de software para el control de sistemas o procesos.
% VI.  Puesta en operación, optimización o automatización de un proceso o sistema industrial.
\newcommand{\lineaTrabajo}{Línea de investigación: IV. Diseño e implementación de sistemas o técnicas de control.}

% Fecha de presentación y defensa oral (Ejemplo: "enero de 2026"):
\newcommand{\fechaDefensa}{enero de 2026}

% ------------------------------------------------------------------------------
% 3. INTEGRANTES DEL EQUIPO (ALUMNOS)
% ------------------------------------------------------------------------------
% La plantilla soporta automáticamente 1, 2 o 3 alumnos.
% Si tu equipo es de 1 o 2 integrantes, deja vacíos los campos no utilizados.

% Alumno 1 (Obligatorio)
\newcommand{\alumnoANombre}{Nombre Completo del Alumno 1}
\newcommand{\alumnoABoleta}{2020XXXXXX}

% Alumno 2 (Opcional - dejar vacío si es individual)
\newcommand{\alumnoBNombre}{Nombre Completo del Alumno 2}
\newcommand{\alumnoBBoleta}{2020XXXXXX}

% Alumno 3 (Opcional - dejar vacío si son 2 o individual)
\newcommand{\alumnoCNombre}{}
\newcommand{\alumnoCBoleta}{}

% ------------------------------------------------------------------------------
% 4. ASESORES DEL PROYECTO
% ------------------------------------------------------------------------------
% La plantilla soporta automáticamente 1, 2 o 3 asesores.
% Si tienes menos de 3 asesores, deja vacíos los campos no utilizados.

% Asesor 1 (Obligatorio)
\newcommand{\asesorANombre}{M. en C. Nombre del Primer Asesor}

% Asesor 2 (Opcional - dejar vacío si no aplica)
\newcommand{\asesorBNombre}{Dra. Nombre del Segundo Asesor}

% Asesor 3 (Opcional - dejar vacío si no aplica)
\newcommand{\asesorCNombre}{}

% ------------------------------------------------------------------------------
% 5. JURADO CALIFICADOR / SÍNODO (PORTADA INTERNA)
% ------------------------------------------------------------------------------
% Estos datos se imprimen en la portada interna con sus líneas de firma.

\newcommand{\juradoPresidente}{Dr. Nombre del Presidente del Jurado}
\newcommand{\juradoSecretario}{M. en C. Nombre del Secretario del Jurado}
\newcommand{\juradoPrimerVocal}{Dr. Nombre del Primer Vocal}
\newcommand{\juradoSegundoVocal}{M. en C. Nombre del Segundo Vocal}
\newcommand{\juradoTercerVocal}{Ing. Nombre del Tercer Vocal (Suplente)}

% ------------------------------------------------------------------------------
% 6. DATOS INSTITUCIONALES (Generalmente no requieren edición)
% ------------------------------------------------------------------------------
\newcommand{\institucionNombre}{INSTITUTO POLITÉCNICO NACIONAL}
\newcommand{\unidadAcademica}{UNIDAD PROFESIONAL INTERDISCIPLINARIA DE INGENIERÍA CAMPUS ZACATECAS}
\newcommand{\siglasUnidad}{UPIIZ}
\newcommand{\programaAcademico}{INGENIERÍA MECATRÓNICA}
\newcommand{\tipoTrabajo}{TRABAJO TERMINAL}
\newcommand{\gradoObtener}{Ingeniero en Mecatrónica}
\newcommand{\lugarDefensa}{Zacatecas, Zac.}
